\documentclass[11pt,a4paper]{article}
\usepackage[utf8]{inputenc}
\usepackage[T1]{fontenc}
\usepackage[spanish,es-minimal,es-noshorthands]{babel}
\AtBeginDocument{\renewcommand{\sin}{\operatorname{sen}}\renewcommand{\arcsin}{\operatorname{arcsen}}}
\usepackage{lmodern}
\usepackage[sfdefault,lining]{FiraSans}
\usepackage[a4paper,top=2.3cm,bottom=2.2cm,left=1.8cm,right=1.8cm,headheight=1.3cm,headsep=0.4cm,footskip=1cm]{geometry}
\usepackage{amsmath,amssymb}
\usepackage{icomma}
\usepackage{graphicx}
\usepackage[table]{xcolor}
\usepackage{tikz}
\usepackage{fontawesome5}
\usepackage{ccicons}
\usepackage{fancyhdr}
\usepackage{lastpage}
\usepackage{tasks}
\usepackage{needspace}
\usepackage[most]{tcolorbox}
\usepackage{array}
\usepackage[hidelinks]{hyperref}
\AtBeginDocument{\hypersetup{pdftitle={Ficha de ejercicios de la unidad 9: Gráficas de funciones. 1º ESO},pdfsubject={Matemáticas 1º ESO},pdfkeywords={coordenadas cartesianas, ejes, cuadrantes, abscisa, ordenada, planos, desplazamientos, simetrías, tablas de valores, gráficas, funciones, crecimiento, máximos y mínimos, proporcionalidad directa, pendiente, función lineal, ordenada en el origen}}}

\definecolor{navy}{HTML}{1B2748}
\definecolor{accent}{HTML}{E8603A}
\definecolor{teal}{HTML}{1F9E90}
\definecolor{soft}{HTML}{EEF1F8}
\definecolor{muted}{HTML}{5D6685}

\setlength{\parindent}{0pt}
\setlength{\parskip}{3pt}
\settasks{label=\alph*), label-format=\bfseries\color{navy}, label-width=1.3em, item-indent=1.8em, column-sep=0.6em, after-item-skip=5pt, before-skip=3pt, after-skip=4pt}

% ---------- capçalera i peu ----------
\newcommand{\llicencia}{\href{https://creativecommons.org/licenses/by-nc-sa/4.0/deed.es}{\ccbyncsa}\ \ Licencia Creative Commons Atribución-NoComercial-CompartirIgual 4.0 Internacional (CC BY-NC-SA 4.0)}
\pagestyle{fancy}
\fancyhf{}
\fancyhead[L]{\raisebox{-0.35\height}{\includegraphics[height=0.95cm]{logo-gva.png}}}
\fancyhead[C]{\small\color{muted}Ficha de ejercicios · Unidad 9 · 1\textsuperscript{o} ESO}
\fancyhead[R]{\raisebox{-0.35\height}{\includegraphics[height=0.95cm]{logo-vermellar.png}}}
\fancyfoot[L]{\scriptsize\color{muted}\llicencia}
\fancyfoot[R]{\small\color{muted}\thepage}
\renewcommand{\headrulewidth}{0.4pt}
\renewcommand{\headrule}{\color{muted!40}\hrule height \headrulewidth}
\fancypagestyle{primera}{\fancyhf{}\renewcommand{\headrulewidth}{0pt}
  \fancyfoot[L]{\scriptsize\color{muted}\llicencia}
  \fancyfoot[R]{\small\color{muted}\thepage}}

% ---------- nivells i exercicis ----------
\newcommand{\estrelles}[1]{{\color{accent}\ifcase#1\or\faStar\faStar[regular]\faStar[regular]\or\faStar\faStar\faStar[regular]\or\faStar\faStar\faStar\fi}}
\newcounter{exer}
\newcommand{\ex}[1]{\par\needspace{4\baselineskip}\medskip\stepcounter{exer}\noindent{\color{navy}\bfseries\theexer.}\ \ {\small\estrelles{#1}}\ \ }

\newcommand{\seccio}[2]{\par\needspace{10\baselineskip}\bigskip
  \begin{tcolorbox}[colback=navy,colframe=navy,coltext=white,arc=2mm,boxrule=0pt,left=3mm,right=3mm,top=1.2mm,bottom=1.2mm]
  {\large\bfseries #1\quad #2}
  \end{tcolorbox}\vspace{-2pt}}

\newcolumntype{C}[1]{>{\centering\arraybackslash}p{#1}}
\definecolor{violet}{HTML}{7C6CF0}

% quadre de recordatori (només a la fitxa d'exercicis)
\newcommand{\recorda}[1]{\par\smallskip\begin{tcolorbox}[colback=accent!7,colframe=accent!55,arc=2mm,boxrule=0.5pt,left=2.5mm,right=2.5mm,top=1mm,bottom=1mm]\small{\color{accent}\faLightbulb[regular]}\ \textbf{Recuerda:}\ #1\end{tcolorbox}}
\newcommand{\res}[1]{#1}\newcommand{\rest}[1]{#1}

\def\fmneed{14}
% =====================================================================
% Capçalera FentMates de les fitxes: un color per a cada curs.
% Els colors són els de la pàgina inicial (LEVELS a assets/js/materials.js);
% el logo ve de _eines/logo/pdf/ (python3 _eines/logo/genera.py).
%
% Ús, al final del preàmbul de la unitat:
%   % =====================================================================
% Capçalera FentMates de les fitxes: un color per a cada curs.
% Els colors són els de la pàgina inicial (LEVELS a assets/js/materials.js);
% el logo ve de _eines/logo/pdf/ (python3 _eines/logo/genera.py).
%
% Ús, al final del preàmbul de la unitat:
%   % =====================================================================
% Capçalera FentMates de les fitxes: un color per a cada curs.
% Els colors són els de la pàgina inicial (LEVELS a assets/js/materials.js);
% el logo ve de _eines/logo/pdf/ (python3 _eines/logo/genera.py).
%
% Ús, al final del preàmbul de la unitat:
%   \input{../comu/fentmates.tex}
%   \fmcurs{1eso}
% i, al començament del document:
%   \capcalera{Matemàtiques}{1r ESO}{Fitxa d'exercicis · Unitat 2}{Divisibilitat}
% Opcional abans de \input: \def\fmneed{14} (espai mínim abans de cada secció).
% =====================================================================
\makeatletter
\def\fm@logodir{../../logo/pdf}
\providecommand{\fmneed}{10}
\def\fm@cursnom{}

% clau del curs, color i color del text damunt del color
\newcommand{\fm@curs}[3]{\@namedef{fm@color@#1}{#2}\@namedef{fm@text@#1}{#3}}
\fm@curs{1eso}{E6B84A}{navy}
\fm@curs{2eso}{3B82F6}{white}
\fm@curs{3eso}{FF7A50}{white}
\fm@curs{4eso-a}{A78BFA}{white}
\fm@curs{4eso-b}{6D28D9}{white}
\fm@curs{1batx-mi}{9B1C1C}{white}
\fm@curs{1batx-ccss}{EF6B6B}{white}
\fm@curs{2batx-mii}{1E6B47}{white}
\fm@curs{2batx-ccss}{2BB3A3}{white}

\newcommand{\fmcurs}[1]{%
  \@ifundefined{fm@color@#1}{\PackageError{fentmates}{Curs desconegut: #1}{}}{}%
  \edef\fm@tmp{\noexpand\definecolor{curs}{HTML}{\@nameuse{fm@color@#1}}%
               \noexpand\colorlet{curstext}{\@nameuse{fm@text@#1}}}\fm@tmp
  \colorlet{cursfosc}{curs!72!black}%
  \def\fm@logo{\fm@logodir/vertical-#1.pdf}%
  \def\fm@marca{\fm@logodir/marca-#1.pdf}}

% ---------- capçalera de la primera pàgina ----------
% #1 matèria (petit), #2 curs, #3 tipus de document i unitat, #4 títol
\newcommand{\capcalera}[4]{%
  \gdef\fm@cursnom{#2}%
  \par\noindent
  \begin{tcolorbox}[colback=curs,colframe=curs,boxrule=0pt,arc=2.2mm,width=3.5cm,height=2.15cm,
      valign=center,left=3mm,right=2mm,top=1mm,bottom=1mm,nobeforeafter,box align=center]
    \color{curstext}{\footnotesize #1}\par\vspace{1pt}{\fontsize{19}{21}\selectfont\bfseries #2\par}%
  \end{tcolorbox}\hspace{5mm}%
  \begin{minipage}[c]{\dimexpr\linewidth-3.5cm-5mm-2.9cm\relax}
    {\color{cursfosc}\small\bfseries\MakeUppercase{#3}}\par\vspace{2pt}
    {\color{navy}\fontsize{20}{23}\selectfont\bfseries #4\par}
  \end{minipage}\hfill
  \begin{minipage}[c]{2.6cm}\raggedleft\includegraphics[width=2.6cm]{\fm@logo}\end{minipage}\par
  \vspace{2pt}}

% ---------- barres de secció del color del curs ----------
\renewcommand{\seccio}[2]{\par\needspace{\fmneed\baselineskip}\bigskip
  \begin{tcolorbox}[colback=curs,colframe=curs,coltext=curstext,arc=2mm,boxrule=0pt,left=3mm,right=3mm,top=1.2mm,bottom=1.2mm]
  {\large\bfseries #1\quad #2}
  \end{tcolorbox}\vspace{-2pt}}

% ---------- estrelles de dificultat del color del curs (a Batxillerat no se'n posen) ----------
\@ifundefined{estrelles}{}{%
  \renewcommand{\estrelles}[1]{{\color{curs}\ifcase#1\or\faStar\faStar[regular]\faStar[regular]\or\faStar\faStar\faStar[regular]\or\faStar\faStar\faStar\fi}}}

% ---------- capçalera i peu de les pàgines ----------
% primera pàgina: logos institucionals al peu
\fancypagestyle{primera}{\fancyhf{}\renewcommand{\headrulewidth}{0pt}
  \fancyfoot[L]{\raisebox{-0.32\height}{\includegraphics[height=0.78cm]{logo-gva.png}}\hspace{5mm}%
    \raisebox{-0.32\height}{\includegraphics[height=0.78cm]{logo-vermellar.png}}}
  \fancyfoot[R]{\scriptsize\color{muted}\href{https://creativecommons.org/licenses/by-nc-sa/4.0/deed.ca}{\ccbyncsa}\ \ CC BY-NC-SA 4.0 · fentmates.org\qquad\small\thepage}}
% la resta: el símbol a l'esquerra i el curs a la dreta, amb el títol al centre com abans
\fancyhead[L]{\raisebox{-0.3\height}{\includegraphics[height=0.62cm]{\fm@marca}}}
\fancyhead[R]{\small\bfseries\color{cursfosc}\fm@cursnom}
\renewcommand{\headrule}{\color{curs!45}\hrule height \headrulewidth}
\makeatother

%   \fmcurs{1eso}
% i, al començament del document:
%   \capcalera{Matemàtiques}{1r ESO}{Fitxa d'exercicis · Unitat 2}{Divisibilitat}
% Opcional abans de \input: \def\fmneed{14} (espai mínim abans de cada secció).
% =====================================================================
\makeatletter
\def\fm@logodir{../../logo/pdf}
\providecommand{\fmneed}{10}
\def\fm@cursnom{}

% clau del curs, color i color del text damunt del color
\newcommand{\fm@curs}[3]{\@namedef{fm@color@#1}{#2}\@namedef{fm@text@#1}{#3}}
\fm@curs{1eso}{E6B84A}{navy}
\fm@curs{2eso}{3B82F6}{white}
\fm@curs{3eso}{FF7A50}{white}
\fm@curs{4eso-a}{A78BFA}{white}
\fm@curs{4eso-b}{6D28D9}{white}
\fm@curs{1batx-mi}{9B1C1C}{white}
\fm@curs{1batx-ccss}{EF6B6B}{white}
\fm@curs{2batx-mii}{1E6B47}{white}
\fm@curs{2batx-ccss}{2BB3A3}{white}

\newcommand{\fmcurs}[1]{%
  \@ifundefined{fm@color@#1}{\PackageError{fentmates}{Curs desconegut: #1}{}}{}%
  \edef\fm@tmp{\noexpand\definecolor{curs}{HTML}{\@nameuse{fm@color@#1}}%
               \noexpand\colorlet{curstext}{\@nameuse{fm@text@#1}}}\fm@tmp
  \colorlet{cursfosc}{curs!72!black}%
  \def\fm@logo{\fm@logodir/vertical-#1.pdf}%
  \def\fm@marca{\fm@logodir/marca-#1.pdf}}

% ---------- capçalera de la primera pàgina ----------
% #1 matèria (petit), #2 curs, #3 tipus de document i unitat, #4 títol
\newcommand{\capcalera}[4]{%
  \gdef\fm@cursnom{#2}%
  \par\noindent
  \begin{tcolorbox}[colback=curs,colframe=curs,boxrule=0pt,arc=2.2mm,width=3.5cm,height=2.15cm,
      valign=center,left=3mm,right=2mm,top=1mm,bottom=1mm,nobeforeafter,box align=center]
    \color{curstext}{\footnotesize #1}\par\vspace{1pt}{\fontsize{19}{21}\selectfont\bfseries #2\par}%
  \end{tcolorbox}\hspace{5mm}%
  \begin{minipage}[c]{\dimexpr\linewidth-3.5cm-5mm-2.9cm\relax}
    {\color{cursfosc}\small\bfseries\MakeUppercase{#3}}\par\vspace{2pt}
    {\color{navy}\fontsize{20}{23}\selectfont\bfseries #4\par}
  \end{minipage}\hfill
  \begin{minipage}[c]{2.6cm}\raggedleft\includegraphics[width=2.6cm]{\fm@logo}\end{minipage}\par
  \vspace{2pt}}

% ---------- barres de secció del color del curs ----------
\renewcommand{\seccio}[2]{\par\needspace{\fmneed\baselineskip}\bigskip
  \begin{tcolorbox}[colback=curs,colframe=curs,coltext=curstext,arc=2mm,boxrule=0pt,left=3mm,right=3mm,top=1.2mm,bottom=1.2mm]
  {\large\bfseries #1\quad #2}
  \end{tcolorbox}\vspace{-2pt}}

% ---------- estrelles de dificultat del color del curs (a Batxillerat no se'n posen) ----------
\@ifundefined{estrelles}{}{%
  \renewcommand{\estrelles}[1]{{\color{curs}\ifcase#1\or\faStar\faStar[regular]\faStar[regular]\or\faStar\faStar\faStar[regular]\or\faStar\faStar\faStar\fi}}}

% ---------- capçalera i peu de les pàgines ----------
% primera pàgina: logos institucionals al peu
\fancypagestyle{primera}{\fancyhf{}\renewcommand{\headrulewidth}{0pt}
  \fancyfoot[L]{\raisebox{-0.32\height}{\includegraphics[height=0.78cm]{logo-gva.png}}\hspace{5mm}%
    \raisebox{-0.32\height}{\includegraphics[height=0.78cm]{logo-vermellar.png}}}
  \fancyfoot[R]{\scriptsize\color{muted}\href{https://creativecommons.org/licenses/by-nc-sa/4.0/deed.ca}{\ccbyncsa}\ \ CC BY-NC-SA 4.0 · fentmates.org\qquad\small\thepage}}
% la resta: el símbol a l'esquerra i el curs a la dreta, amb el títol al centre com abans
\fancyhead[L]{\raisebox{-0.3\height}{\includegraphics[height=0.62cm]{\fm@marca}}}
\fancyhead[R]{\small\bfseries\color{cursfosc}\fm@cursnom}
\renewcommand{\headrule}{\color{curs!45}\hrule height \headrulewidth}
\makeatother

%   \fmcurs{1eso}
% i, al començament del document:
%   \capcalera{Matemàtiques}{1r ESO}{Fitxa d'exercicis · Unitat 2}{Divisibilitat}
% Opcional abans de \input: \def\fmneed{14} (espai mínim abans de cada secció).
% =====================================================================
\makeatletter
\def\fm@logodir{../../logo/pdf}
\providecommand{\fmneed}{10}
\def\fm@cursnom{}

% clau del curs, color i color del text damunt del color
\newcommand{\fm@curs}[3]{\@namedef{fm@color@#1}{#2}\@namedef{fm@text@#1}{#3}}
\fm@curs{1eso}{E6B84A}{navy}
\fm@curs{2eso}{3B82F6}{white}
\fm@curs{3eso}{FF7A50}{white}
\fm@curs{4eso-a}{A78BFA}{white}
\fm@curs{4eso-b}{6D28D9}{white}
\fm@curs{1batx-mi}{9B1C1C}{white}
\fm@curs{1batx-ccss}{EF6B6B}{white}
\fm@curs{2batx-mii}{1E6B47}{white}
\fm@curs{2batx-ccss}{2BB3A3}{white}

\newcommand{\fmcurs}[1]{%
  \@ifundefined{fm@color@#1}{\PackageError{fentmates}{Curs desconegut: #1}{}}{}%
  \edef\fm@tmp{\noexpand\definecolor{curs}{HTML}{\@nameuse{fm@color@#1}}%
               \noexpand\colorlet{curstext}{\@nameuse{fm@text@#1}}}\fm@tmp
  \colorlet{cursfosc}{curs!72!black}%
  \def\fm@logo{\fm@logodir/vertical-#1.pdf}%
  \def\fm@marca{\fm@logodir/marca-#1.pdf}}

% ---------- capçalera de la primera pàgina ----------
% #1 matèria (petit), #2 curs, #3 tipus de document i unitat, #4 títol
\newcommand{\capcalera}[4]{%
  \gdef\fm@cursnom{#2}%
  \par\noindent
  \begin{tcolorbox}[colback=curs,colframe=curs,boxrule=0pt,arc=2.2mm,width=3.5cm,height=2.15cm,
      valign=center,left=3mm,right=2mm,top=1mm,bottom=1mm,nobeforeafter,box align=center]
    \color{curstext}{\footnotesize #1}\par\vspace{1pt}{\fontsize{19}{21}\selectfont\bfseries #2\par}%
  \end{tcolorbox}\hspace{5mm}%
  \begin{minipage}[c]{\dimexpr\linewidth-3.5cm-5mm-2.9cm\relax}
    {\color{cursfosc}\small\bfseries\MakeUppercase{#3}}\par\vspace{2pt}
    {\color{navy}\fontsize{20}{23}\selectfont\bfseries #4\par}
  \end{minipage}\hfill
  \begin{minipage}[c]{2.6cm}\raggedleft\includegraphics[width=2.6cm]{\fm@logo}\end{minipage}\par
  \vspace{2pt}}

% ---------- barres de secció del color del curs ----------
\renewcommand{\seccio}[2]{\par\needspace{\fmneed\baselineskip}\bigskip
  \begin{tcolorbox}[colback=curs,colframe=curs,coltext=curstext,arc=2mm,boxrule=0pt,left=3mm,right=3mm,top=1.2mm,bottom=1.2mm]
  {\large\bfseries #1\quad #2}
  \end{tcolorbox}\vspace{-2pt}}

% ---------- estrelles de dificultat del color del curs (a Batxillerat no se'n posen) ----------
\@ifundefined{estrelles}{}{%
  \renewcommand{\estrelles}[1]{{\color{curs}\ifcase#1\or\faStar\faStar[regular]\faStar[regular]\or\faStar\faStar\faStar[regular]\or\faStar\faStar\faStar\fi}}}

% ---------- capçalera i peu de les pàgines ----------
% primera pàgina: logos institucionals al peu
\fancypagestyle{primera}{\fancyhf{}\renewcommand{\headrulewidth}{0pt}
  \fancyfoot[L]{\raisebox{-0.32\height}{\includegraphics[height=0.78cm]{logo-gva.png}}\hspace{5mm}%
    \raisebox{-0.32\height}{\includegraphics[height=0.78cm]{logo-vermellar.png}}}
  \fancyfoot[R]{\scriptsize\color{muted}\href{https://creativecommons.org/licenses/by-nc-sa/4.0/deed.ca}{\ccbyncsa}\ \ CC BY-NC-SA 4.0 · fentmates.org\qquad\small\thepage}}
% la resta: el símbol a l'esquerra i el curs a la dreta, amb el títol al centre com abans
\fancyhead[L]{\raisebox{-0.3\height}{\includegraphics[height=0.62cm]{\fm@marca}}}
\fancyhead[R]{\small\bfseries\color{cursfosc}\fm@cursnom}
\renewcommand{\headrule}{\color{curs!45}\hrule height \headrulewidth}
\makeatother

\fmcurs{1eso}

\begin{document}
\thispagestyle{primera}

\capcalera{Matemáticas}{1\textsuperscript{o} ESO}{Ficha de ejercicios · Unidad 9}{Gráficas de funciones}
\vspace{10pt}
Nombre y apellidos: \rule{6.6cm}{0.4pt}\hfill Grupo: \rule{1.6cm}{0.4pt}\hfill Fecha: \rule{2.2cm}{0.4pt}
\vspace{8pt}

\begin{tcolorbox}[colback=soft,colframe=soft,arc=2mm,left=3mm,right=3mm,top=1.5mm,bottom=1.5mm]
\small Ejercicios graduados por niveles de dificultad:\quad \estrelles{1}\ básico \qquad \estrelles{2}\ medio \qquad \estrelles{3}\ avanzado.\\[2pt] Resuélvelos en el cuaderno mostrando todos los pasos. Para dibujar puntos, rectas y gráficas usa papel cuadriculado y escribe siempre el nombre y la unidad de cada eje.
\end{tcolorbox}

% =====================================================
\seccio{1}{El plano cartesiano. Coordenadas de un punto}

\ex{1} Escribe las coordenadas de los puntos del dibujo:\par\smallskip\begin{center}\begin{tikzpicture}[x=0.42cm,y=0.42cm,baseline=(current bounding box.center),line cap=round,line join=round]
\draw[muted!30,line width=0.3pt] (-5,-5)--(-5,5);
\draw[muted!30,line width=0.3pt] (-4,-5)--(-4,5);
\draw[muted!30,line width=0.3pt] (-3,-5)--(-3,5);
\draw[muted!30,line width=0.3pt] (-2,-5)--(-2,5);
\draw[muted!30,line width=0.3pt] (-1,-5)--(-1,5);
\draw[muted!30,line width=0.3pt] (0,-5)--(0,5);
\draw[muted!30,line width=0.3pt] (1,-5)--(1,5);
\draw[muted!30,line width=0.3pt] (2,-5)--(2,5);
\draw[muted!30,line width=0.3pt] (3,-5)--(3,5);
\draw[muted!30,line width=0.3pt] (4,-5)--(4,5);
\draw[muted!30,line width=0.3pt] (5,-5)--(5,5);
\draw[muted!30,line width=0.3pt] (-5,-5)--(5,-5);
\draw[muted!30,line width=0.3pt] (-5,-4)--(5,-4);
\draw[muted!30,line width=0.3pt] (-5,-3)--(5,-3);
\draw[muted!30,line width=0.3pt] (-5,-2)--(5,-2);
\draw[muted!30,line width=0.3pt] (-5,-1)--(5,-1);
\draw[muted!30,line width=0.3pt] (-5,0)--(5,0);
\draw[muted!30,line width=0.3pt] (-5,1)--(5,1);
\draw[muted!30,line width=0.3pt] (-5,2)--(5,2);
\draw[muted!30,line width=0.3pt] (-5,3)--(5,3);
\draw[muted!30,line width=0.3pt] (-5,4)--(5,4);
\draw[muted!30,line width=0.3pt] (-5,5)--(5,5);
\draw[-latex,muted,line width=0.6pt] (-5,0)--(5.3,0) node[right,font=\scriptsize,text=navy]{x};
\draw[-latex,muted,line width=0.6pt] (0,-5)--(0,5.3) node[above,font=\scriptsize,text=navy]{y};
\node[font=\tiny,text=muted,below,inner sep=1.5pt] at (-4,0) {$-4$};
\node[font=\tiny,text=muted,below,inner sep=1.5pt] at (-3,0) {$-3$};
\node[font=\tiny,text=muted,below,inner sep=1.5pt] at (-2,0) {$-2$};
\node[font=\tiny,text=muted,below,inner sep=1.5pt] at (-1,0) {$-1$};
\node[font=\tiny,text=muted,below,inner sep=1.5pt] at (1,0) {$1$};
\node[font=\tiny,text=muted,below,inner sep=1.5pt] at (2,0) {$2$};
\node[font=\tiny,text=muted,below,inner sep=1.5pt] at (3,0) {$3$};
\node[font=\tiny,text=muted,below,inner sep=1.5pt] at (4,0) {$4$};
\node[font=\tiny,text=muted,left,inner sep=1.5pt] at (0,-4) {$-4$};
\node[font=\tiny,text=muted,left,inner sep=1.5pt] at (0,-3) {$-3$};
\node[font=\tiny,text=muted,left,inner sep=1.5pt] at (0,-2) {$-2$};
\node[font=\tiny,text=muted,left,inner sep=1.5pt] at (0,-1) {$-1$};
\node[font=\tiny,text=muted,left,inner sep=1.5pt] at (0,1) {$1$};
\node[font=\tiny,text=muted,left,inner sep=1.5pt] at (0,2) {$2$};
\node[font=\tiny,text=muted,left,inner sep=1.5pt] at (0,3) {$3$};
\node[font=\tiny,text=muted,left,inner sep=1.5pt] at (0,4) {$4$};
\node[font=\tiny,text=muted,below left,inner sep=1.5pt] at (0,0) {$0$};
\fill[navy] (2,3) circle (2.2pt) node[above right,font=\scriptsize,inner sep=1.5pt]{$A$};
\fill[navy] (-4,1) circle (2.2pt) node[above right,font=\scriptsize,inner sep=1.5pt]{$B$};
\fill[navy] (-3,-2) circle (2.2pt) node[above right,font=\scriptsize,inner sep=1.5pt]{$C$};
\fill[navy] (3,-4) circle (2.2pt) node[above right,font=\scriptsize,inner sep=1.5pt]{$D$};
\fill[navy] (0,2) circle (2.2pt) node[above right,font=\scriptsize,inner sep=1.5pt]{$E$};
\fill[navy] (4,0) circle (2.2pt) node[above right,font=\scriptsize,inner sep=1.5pt]{$F$};
\fill[navy] (-1,4) circle (2.2pt) node[above right,font=\scriptsize,inner sep=1.5pt]{$G$};
\fill[navy] (-2,0) circle (2.2pt) node[above right,font=\scriptsize,inner sep=1.5pt]{$H$};
\end{tikzpicture}\end{center}\recorda{Las coordenadas de un punto se escriben $P(x,\ y)$: primero la \textbf{abscisa} $x$ (horizontal: derecha $+$, izquierda $-$) y después la \textbf{ordenada} $y$ (vertical: arriba $+$, abajo $-$).}

\ex{1} Dibuja en unos ejes de coordenadas los puntos $P(3,\ 1)$, $Q(-2,\ 4)$, $R(-4,\ -3)$, $S(1,\ -2)$, $T(0,\ -4)$ y $U(5,\ 0)$.

\ex{1} Sin dibujarlos, di en qué cuadrante o sobre qué eje está cada punto:\begin{tasks}(4)
  \task $(6,\ -2)$
  \task $(-3,\ -8)$
  \task $(-1,\ 7)$
  \task $(0,\ -5)$
  \task $(9,\ 0)$
  \task $(4,\ 4)$
  \task $(-6,\ 0)$
  \task $(-2,\ -2)$
\end{tasks}\recorda{Signos: cuadrante I $(+,\,+)$, II $(-,\,+)$, III $(-,\,-)$, IV $(+,\,-)$. Si $y=0$, el punto es del eje $X$; si $x=0$, del eje $Y$.}

\ex{1} Escribe las coordenadas del punto:\begin{tasks}(1)
  \task que está 3 unidades a la derecha del eje $Y$ y 5 por encima del eje $X$;
  \task que está 2 unidades a la izquierda del eje $Y$ y 4 por debajo del eje $X$;
  \task que está sobre el eje $X$, 6 unidades a la izquierda del origen;
  \task que está sobre el eje $Y$, 1 unidad por encima del origen.
\end{tasks}

\ex{2} Dibuja los puntos $(-2,\ 4)$ y $(4,\ -2)$. ¿Son el mismo punto? ¿En qué cuadrante está cada uno?

\ex{2} Responde:\begin{tasks}(1)
  \task ¿Qué punto del cuadrante II está a 4 unidades de cada eje?
  \task ¿Qué puntos del eje $X$ están a 3 unidades del origen?
  \task ¿Qué punto del cuadrante III tiene la abscisa $-5$ y la ordenada la mitad de la abscisa?
\end{tasks}

\ex{2} El punto $(a,\ b)$ está en el cuadrante IV. ¿En qué cuadrante están los puntos $(-a,\ b)$, \ $(a,\ -b)$ \ y \ $(b,\ a)$?

\ex{3} Dibuja todos los puntos de coordenadas enteras entre $-4$ y $4$ que tienen la ordenada igual a la abscisa. Haz lo mismo con los que tienen la ordenada igual a 2. ¿Qué observas?

\ex{3} Dibuja los puntos $A(-3,\ -2)$, $B(1,\ -2)$ y $C(1,\ 2)$. Halla el punto $D$ para que $ABCD$ sea un cuadrado y calcula las coordenadas del centro del cuadrado.


% =====================================================
\seccio{2}{Aplicaciones de las coordenadas}

\ex{1} En este plano de un pueblo, la bandera es la plaza (el origen) y cada cuadro es una manzana de casas. Escribe las coordenadas de la escuela ($E$), el parque ($P$), el hospital ($H$) y el mercado ($M$). ¿Qué hay en $(0,\ -1)$? ¿Y en $(-4,\ 0)$?\par\smallskip\begin{center}\begin{tikzpicture}[x=0.42cm,y=0.42cm,baseline=(current bounding box.center),line cap=round,line join=round]
\draw[muted!30,line width=0.3pt] (-5,-5)--(-5,5);
\draw[muted!30,line width=0.3pt] (-4,-5)--(-4,5);
\draw[muted!30,line width=0.3pt] (-3,-5)--(-3,5);
\draw[muted!30,line width=0.3pt] (-2,-5)--(-2,5);
\draw[muted!30,line width=0.3pt] (-1,-5)--(-1,5);
\draw[muted!30,line width=0.3pt] (0,-5)--(0,5);
\draw[muted!30,line width=0.3pt] (1,-5)--(1,5);
\draw[muted!30,line width=0.3pt] (2,-5)--(2,5);
\draw[muted!30,line width=0.3pt] (3,-5)--(3,5);
\draw[muted!30,line width=0.3pt] (4,-5)--(4,5);
\draw[muted!30,line width=0.3pt] (5,-5)--(5,5);
\draw[muted!30,line width=0.3pt] (-5,-5)--(5,-5);
\draw[muted!30,line width=0.3pt] (-5,-4)--(5,-4);
\draw[muted!30,line width=0.3pt] (-5,-3)--(5,-3);
\draw[muted!30,line width=0.3pt] (-5,-2)--(5,-2);
\draw[muted!30,line width=0.3pt] (-5,-1)--(5,-1);
\draw[muted!30,line width=0.3pt] (-5,0)--(5,0);
\draw[muted!30,line width=0.3pt] (-5,1)--(5,1);
\draw[muted!30,line width=0.3pt] (-5,2)--(5,2);
\draw[muted!30,line width=0.3pt] (-5,3)--(5,3);
\draw[muted!30,line width=0.3pt] (-5,4)--(5,4);
\draw[muted!30,line width=0.3pt] (-5,5)--(5,5);
\draw[-latex,muted,line width=0.6pt] (-5,0)--(5.3,0) node[right,font=\scriptsize,text=navy]{x};
\draw[-latex,muted,line width=0.6pt] (0,-5)--(0,5.3) node[above,font=\scriptsize,text=navy]{y};
\node[font=\tiny,text=muted,below,inner sep=1.5pt] at (-4,0) {$-4$};
\node[font=\tiny,text=muted,below,inner sep=1.5pt] at (-3,0) {$-3$};
\node[font=\tiny,text=muted,below,inner sep=1.5pt] at (-2,0) {$-2$};
\node[font=\tiny,text=muted,below,inner sep=1.5pt] at (-1,0) {$-1$};
\node[font=\tiny,text=muted,below,inner sep=1.5pt] at (1,0) {$1$};
\node[font=\tiny,text=muted,below,inner sep=1.5pt] at (2,0) {$2$};
\node[font=\tiny,text=muted,below,inner sep=1.5pt] at (3,0) {$3$};
\node[font=\tiny,text=muted,below,inner sep=1.5pt] at (4,0) {$4$};
\node[font=\tiny,text=muted,left,inner sep=1.5pt] at (0,-4) {$-4$};
\node[font=\tiny,text=muted,left,inner sep=1.5pt] at (0,-3) {$-3$};
\node[font=\tiny,text=muted,left,inner sep=1.5pt] at (0,-2) {$-2$};
\node[font=\tiny,text=muted,left,inner sep=1.5pt] at (0,-1) {$-1$};
\node[font=\tiny,text=muted,left,inner sep=1.5pt] at (0,1) {$1$};
\node[font=\tiny,text=muted,left,inner sep=1.5pt] at (0,2) {$2$};
\node[font=\tiny,text=muted,left,inner sep=1.5pt] at (0,3) {$3$};
\node[font=\tiny,text=muted,left,inner sep=1.5pt] at (0,4) {$4$};
\node[font=\tiny,text=muted,below left,inner sep=1.5pt] at (0,0) {$0$};
\fill[accent] (2,3) circle (2.2pt) node[above right,font=\scriptsize,inner sep=1.5pt]{\textbf{E}};
\fill[accent] (-3,2) circle (2.2pt) node[above right,font=\scriptsize,inner sep=1.5pt]{\textbf{P}};
\fill[accent] (4,-2) circle (2.2pt) node[above right,font=\scriptsize,inner sep=1.5pt]{\textbf{H}};
\fill[accent] (-2,-3) circle (2.2pt) node[above right,font=\scriptsize,inner sep=1.5pt]{\textbf{M}};
\fill[accent] (0,-1) circle (2.2pt) node[above right,font=\scriptsize,inner sep=1.5pt]{\textbf{B}};
\fill[accent] (-4,0) circle (2.2pt) node[above right,font=\scriptsize,inner sep=1.5pt]{\textbf{C}};
\node[font=\scriptsize,above right,text=muted] at (0.25,0.25) {\faFlag};
\end{tikzpicture}\end{center}

\ex{1} Dibuja el polígono de vértices $A(1,\ 1)$, $B(6,\ 1)$, $C(6,\ 4)$ y $D(1,\ 4)$. ¿Qué polígono es? Calcula su perímetro y su área (la unidad es el cm).\recorda{Si dos puntos tienen la misma $y$, la distancia entre ellos es la diferencia de las $x$; si tienen la misma $x$, la diferencia de las $y$.}

\ex{1} Calcula la distancia entre los puntos:\begin{tasks}(2)
  \task $(-2,\ 3)$ y $(5,\ 3)$
  \task $(4,\ -1)$ y $(4,\ 6)$
  \task $(-6,\ -2)$ y $(-1,\ -2)$
  \task $(0,\ -4)$ y $(0,\ 3)$
\end{tasks}

\ex{2} Un robot sale del origen y hace estos movimientos: 4 a la derecha, 3 arriba, 6 a la izquierda, 4 abajo y 3 a la derecha. ¿Dónde acaba? Dibuja su camino.

\ex{2} Desplaza el triángulo de vértices $(-4,\ 1)$, $(-1,\ 1)$ y $(-2,\ 4)$ 5 unidades a la derecha y 4 hacia abajo. Escribe los vértices nuevos y dibuja los dos triángulos.\recorda{Desplazar $a$ a la derecha y $b$ arriba: $(x,\ y)\to(x+a,\ y+b)$ (si va a la izquierda o abajo, restamos).}

\ex{2} Completa la tabla con los simétricos de cada punto respecto a los ejes:\par\smallskip\begin{center}\renewcommand{\arraystretch}{1.3}\begin{tabular}{|c|C{2.2cm}|C{2.2cm}|}\hline Punto & Simétrico respecto al eje $X$ & Simétrico respecto al eje $Y$\\\hline $(3,\ 5)$ & & \\\hline $(-2,\ 4)$ & & \\\hline $(-1,\ -6)$ & & \\\hline $(5,\ -3)$ & & \\\hline\end{tabular}\end{center}\recorda{Respecto al eje $X$ cambia el signo de la $y$: $(x,\ -y)$. Respecto al eje $Y$ cambia el signo de la $x$: $(-x,\ y)$.}

\ex{2} Tres vértices de un rectángulo son $A(-2,\ 1)$, $B(3,\ 1)$ y $C(3,\ -2)$. Halla el cuarto vértice $D$ y calcula su área y su perímetro.

\ex{3} Dibuja el triángulo de vértices $A(-2,\ -1)$, $B(4,\ -1)$ y $C(1,\ 3)$ y calcula su área.

\ex{3} Dibuja el pentágono de vértices $(0,\ 0)$, $(4,\ 0)$, $(4,\ 2)$, $(2,\ 4)$ y $(0,\ 2)$ y calcula su área.


% =====================================================
\seccio{3}{Relaciones entre magnitudes: tablas y gráficas}

\ex{1} Un panecillo vale 0,40\,€. Completa la tabla, dibuja la gráfica y di si se tienen que unir los puntos.\par\smallskip\begin{center}$\begin{array}{|c|C{0.75cm}C{0.75cm}C{0.75cm}C{0.75cm}C{0.75cm}C{0.75cm}|}\hline \text{Panecillos}&1&2&3&4&5&6\\\hline \text{Precio (€)}&\phantom{00}&\phantom{00}&\phantom{00}&\phantom{00}&\phantom{00}&\phantom{00}\\\hline\end{array}$\end{center}\recorda{Cada par $(x,\ y)$ de la tabla es un punto de la gráfica. Unimos los puntos solo si también tienen sentido los valores intermedios.}

\ex{1} Di cuál es la variable independiente y cuál la dependiente:\begin{tasks}(1)
  \task el tiempo que está abierto un grifo y los litros de agua que salen;
  \task el número de fotocopias y el precio que pagamos;
  \task la hora del día y la temperatura;
  \task los km que hace un coche y la gasolina que le queda.
\end{tasks}

\ex{1} Completa la tabla de la función $y=2x+3$, escribe los puntos y dibújalos.\par\smallskip\begin{center}$\begin{array}{|c|C{0.75cm}C{0.75cm}C{0.75cm}C{0.75cm}C{0.75cm}|}\hline x&0&1&2&3&4\\\hline y&\phantom{00}&\phantom{00}&\phantom{00}&\phantom{00}&\phantom{00}\\\hline\end{array}$\end{center}

\ex{1} La tabla da la temperatura de un día de primavera. Dibuja la gráfica (uniendo los puntos) y di a qué hora hizo más calor y a qué hora más frío.\par\smallskip\begin{center}$\begin{array}{|c|ccccccccc|}\hline \text{Hora}&0&3&6&9&12&15&18&21&24\\\hline ^\circ\text{C}&8&6&7&12&16&18&15&11&9\\\hline\end{array}$\end{center}

\ex{2} ¿Cuáles de estas tablas corresponden a una función? Explica por qué.\par\smallskip\begin{center}$\begin{array}{|c|cccc|}\hline x&1&2&3&4\\\hline y&5&10&15&20\\\hline\end{array}$\quad $\begin{array}{|c|cccc|}\hline x&1&1&2&3\\\hline y&4&6&8&9\\\hline\end{array}$\quad $\begin{array}{|c|cccc|}\hline x&0&2&4&6\\\hline y&3&3&3&3\\\hline\end{array}$\end{center}

\ex{2} Halla la fórmula de cada tabla:\par\smallskip\begin{center}$\begin{array}{|c|cccc|}\hline x&0&1&2&3\\\hline y&5&7&9&11\\\hline\end{array}$\quad $\begin{array}{|c|cccc|}\hline x&1&2&3&4\\\hline y&10&20&30&40\\\hline\end{array}$\quad $\begin{array}{|c|cccc|}\hline x&0&1&2&3\\\hline y&9&8&7&6\\\hline\end{array}$\end{center}

\ex{2} Una tarifa de móvil cuesta 10\,€ al mes más 3\,€ por cada GB de datos. Escribe la fórmula del precio $y$ si gastamos $x$ GB y haz una tabla para 0, 1, 2, 5 y 10 GB.

\ex{3} Haz una tabla del perímetro $P$ y otra del área $A$ de un cuadrado de lado $l=1,\ 2,\ 3,\ 4,\ 5$\,cm. Escribe sus fórmulas y dibuja las gráficas. ¿Cuál es una recta?

\ex{3} Queremos hacer rectángulos de 20\,cm de perímetro con lados de cm enteros. Haz una tabla con la base $x$ y la altura $y$, escribe su fórmula y dibuja la gráfica. ¿Es creciente o decreciente?


% =====================================================
\seccio{4}{Interpretación de gráficas}

\ex{1} La gráfica muestra la excursión en bici de Nuria (distancia a casa a lo largo del tiempo).\par\smallskip\begin{center}\begin{tikzpicture}[x=0.75cm,y=0.3cm,baseline=(current bounding box.center),line cap=round,line join=round]
\draw[muted!30,line width=0.3pt] (0,0)--(0,16);
\draw[muted!30,line width=0.3pt] (1,0)--(1,16);
\draw[muted!30,line width=0.3pt] (2,0)--(2,16);
\draw[muted!30,line width=0.3pt] (3,0)--(3,16);
\draw[muted!30,line width=0.3pt] (4,0)--(4,16);
\draw[muted!30,line width=0.3pt] (5,0)--(5,16);
\draw[muted!30,line width=0.3pt] (6,0)--(6,16);
\draw[muted!30,line width=0.3pt] (7,0)--(7,16);
\draw[muted!30,line width=0.3pt] (0,0)--(7.5,0);
\draw[muted!30,line width=0.3pt] (0,2)--(7.5,2);
\draw[muted!30,line width=0.3pt] (0,4)--(7.5,4);
\draw[muted!30,line width=0.3pt] (0,6)--(7.5,6);
\draw[muted!30,line width=0.3pt] (0,8)--(7.5,8);
\draw[muted!30,line width=0.3pt] (0,10)--(7.5,10);
\draw[muted!30,line width=0.3pt] (0,12)--(7.5,12);
\draw[muted!30,line width=0.3pt] (0,14)--(7.5,14);
\draw[muted!30,line width=0.3pt] (0,16)--(7.5,16);
\draw[-latex,muted,line width=0.6pt] (0,0)--(7.8,0) node[right,font=\scriptsize,text=navy]{tiempo (h)};
\draw[-latex,muted,line width=0.6pt] (0,0)--(0,16.6) node[above,font=\scriptsize,text=navy]{distancia (km)};
\node[font=\tiny,text=muted,below,inner sep=1.5pt] at (1,0) {$1$};
\node[font=\tiny,text=muted,below,inner sep=1.5pt] at (2,0) {$2$};
\node[font=\tiny,text=muted,below,inner sep=1.5pt] at (3,0) {$3$};
\node[font=\tiny,text=muted,below,inner sep=1.5pt] at (4,0) {$4$};
\node[font=\tiny,text=muted,below,inner sep=1.5pt] at (5,0) {$5$};
\node[font=\tiny,text=muted,below,inner sep=1.5pt] at (6,0) {$6$};
\node[font=\tiny,text=muted,below,inner sep=1.5pt] at (7,0) {$7$};
\node[font=\tiny,text=muted,left,inner sep=1.5pt] at (0,2) {$2$};
\node[font=\tiny,text=muted,left,inner sep=1.5pt] at (0,4) {$4$};
\node[font=\tiny,text=muted,left,inner sep=1.5pt] at (0,6) {$6$};
\node[font=\tiny,text=muted,left,inner sep=1.5pt] at (0,8) {$8$};
\node[font=\tiny,text=muted,left,inner sep=1.5pt] at (0,10) {$10$};
\node[font=\tiny,text=muted,left,inner sep=1.5pt] at (0,12) {$12$};
\node[font=\tiny,text=muted,left,inner sep=1.5pt] at (0,14) {$14$};
\node[font=\tiny,text=muted,below left,inner sep=1.5pt] at (0,0) {$0$};
\draw[accent,line width=1.1pt] (0,0)--(1,6)--(2,6)--(4,14)--(5,14)--(7,0);
\fill[navy] (0,0) circle (2.2pt);
\fill[navy] (1,6) circle (2.2pt);
\fill[navy] (2,6) circle (2.2pt);
\fill[navy] (4,14) circle (2.2pt);
\fill[navy] (5,14) circle (2.2pt);
\fill[navy] (7,0) circle (2.2pt);
\end{tikzpicture}\end{center}\begin{tasks}(1)
  \task ¿A qué distancia de casa estaba a las 3 horas? ¿Y a las 6 horas?
  \task ¿Cuántas veces paró y cuánto tiempo cada vez?
  \task ¿Cuál es la distancia máxima a la que llegó? ¿A qué hora?
  \task ¿Cuánto tardó en volver? ¿Qué distancia recorrió en total?
\end{tasks}\recorda{Gráfica de distancia y tiempo: si sube, se aleja; si es horizontal, está parado; si baja, vuelve. Cuanto más inclinada, más deprisa.}

\ex{1} La gráfica da la temperatura a lo largo de un día.\par\smallskip\begin{center}\begin{tikzpicture}[x=0.26cm,y=0.2cm,baseline=(current bounding box.center),line cap=round,line join=round]
\draw[muted!30,line width=0.3pt] (0,0)--(0,20);
\draw[muted!30,line width=0.3pt] (2,0)--(2,20);
\draw[muted!30,line width=0.3pt] (4,0)--(4,20);
\draw[muted!30,line width=0.3pt] (6,0)--(6,20);
\draw[muted!30,line width=0.3pt] (8,0)--(8,20);
\draw[muted!30,line width=0.3pt] (10,0)--(10,20);
\draw[muted!30,line width=0.3pt] (12,0)--(12,20);
\draw[muted!30,line width=0.3pt] (14,0)--(14,20);
\draw[muted!30,line width=0.3pt] (16,0)--(16,20);
\draw[muted!30,line width=0.3pt] (18,0)--(18,20);
\draw[muted!30,line width=0.3pt] (20,0)--(20,20);
\draw[muted!30,line width=0.3pt] (22,0)--(22,20);
\draw[muted!30,line width=0.3pt] (24,0)--(24,20);
\draw[muted!30,line width=0.3pt] (0,0)--(25,0);
\draw[muted!30,line width=0.3pt] (0,2)--(25,2);
\draw[muted!30,line width=0.3pt] (0,4)--(25,4);
\draw[muted!30,line width=0.3pt] (0,6)--(25,6);
\draw[muted!30,line width=0.3pt] (0,8)--(25,8);
\draw[muted!30,line width=0.3pt] (0,10)--(25,10);
\draw[muted!30,line width=0.3pt] (0,12)--(25,12);
\draw[muted!30,line width=0.3pt] (0,14)--(25,14);
\draw[muted!30,line width=0.3pt] (0,16)--(25,16);
\draw[muted!30,line width=0.3pt] (0,18)--(25,18);
\draw[muted!30,line width=0.3pt] (0,20)--(25,20);
\draw[-latex,muted,line width=0.6pt] (0,0)--(25.6,0) node[right,font=\scriptsize,text=navy]{h};
\draw[-latex,muted,line width=0.6pt] (0,0)--(0,20.6) node[above,font=\scriptsize,text=navy]{$^\circ$C};
\node[font=\tiny,text=muted,below,inner sep=1.5pt] at (2,0) {$2$};
\node[font=\tiny,text=muted,below,inner sep=1.5pt] at (4,0) {$4$};
\node[font=\tiny,text=muted,below,inner sep=1.5pt] at (6,0) {$6$};
\node[font=\tiny,text=muted,below,inner sep=1.5pt] at (8,0) {$8$};
\node[font=\tiny,text=muted,below,inner sep=1.5pt] at (10,0) {$10$};
\node[font=\tiny,text=muted,below,inner sep=1.5pt] at (12,0) {$12$};
\node[font=\tiny,text=muted,below,inner sep=1.5pt] at (14,0) {$14$};
\node[font=\tiny,text=muted,below,inner sep=1.5pt] at (16,0) {$16$};
\node[font=\tiny,text=muted,below,inner sep=1.5pt] at (18,0) {$18$};
\node[font=\tiny,text=muted,below,inner sep=1.5pt] at (20,0) {$20$};
\node[font=\tiny,text=muted,below,inner sep=1.5pt] at (22,0) {$22$};
\node[font=\tiny,text=muted,below,inner sep=1.5pt] at (24,0) {$24$};
\node[font=\tiny,text=muted,left,inner sep=1.5pt] at (0,2) {$2$};
\node[font=\tiny,text=muted,left,inner sep=1.5pt] at (0,4) {$4$};
\node[font=\tiny,text=muted,left,inner sep=1.5pt] at (0,6) {$6$};
\node[font=\tiny,text=muted,left,inner sep=1.5pt] at (0,8) {$8$};
\node[font=\tiny,text=muted,left,inner sep=1.5pt] at (0,10) {$10$};
\node[font=\tiny,text=muted,left,inner sep=1.5pt] at (0,12) {$12$};
\node[font=\tiny,text=muted,left,inner sep=1.5pt] at (0,14) {$14$};
\node[font=\tiny,text=muted,left,inner sep=1.5pt] at (0,16) {$16$};
\node[font=\tiny,text=muted,left,inner sep=1.5pt] at (0,18) {$18$};
\node[font=\tiny,text=muted,below left,inner sep=1.5pt] at (0,0) {$0$};
\draw[accent,line width=1.1pt] (0,6)--(2,4)--(6,4)--(10,12)--(14,18)--(18,14)--(22,8)--(24,7);
\end{tikzpicture}\end{center}\begin{tasks}(1)
  \task ¿Cuál fue la temperatura máxima? ¿Y la mínima?
  \task ¿En qué intervalos de tiempo sube la temperatura? ¿En cuáles baja? ¿Cuándo es constante?
\end{tasks}

\ex{1} La gráfica muestra la altura de una planta durante 6 semanas.\par\smallskip\begin{center}\begin{tikzpicture}[x=0.6cm,y=0.22cm,baseline=(current bounding box.center),line cap=round,line join=round]
\draw[muted!30,line width=0.3pt] (0,0)--(0,16.5);
\draw[muted!30,line width=0.3pt] (1,0)--(1,16.5);
\draw[muted!30,line width=0.3pt] (2,0)--(2,16.5);
\draw[muted!30,line width=0.3pt] (3,0)--(3,16.5);
\draw[muted!30,line width=0.3pt] (4,0)--(4,16.5);
\draw[muted!30,line width=0.3pt] (5,0)--(5,16.5);
\draw[muted!30,line width=0.3pt] (6,0)--(6,16.5);
\draw[muted!30,line width=0.3pt] (0,0)--(6.6,0);
\draw[muted!30,line width=0.3pt] (0,2)--(6.6,2);
\draw[muted!30,line width=0.3pt] (0,4)--(6.6,4);
\draw[muted!30,line width=0.3pt] (0,6)--(6.6,6);
\draw[muted!30,line width=0.3pt] (0,8)--(6.6,8);
\draw[muted!30,line width=0.3pt] (0,10)--(6.6,10);
\draw[muted!30,line width=0.3pt] (0,12)--(6.6,12);
\draw[muted!30,line width=0.3pt] (0,14)--(6.6,14);
\draw[muted!30,line width=0.3pt] (0,16)--(6.6,16);
\draw[-latex,muted,line width=0.6pt] (0,0)--(6.9,0) node[right,font=\scriptsize,text=navy]{semanas};
\draw[-latex,muted,line width=0.6pt] (0,0)--(0,17.1) node[above,font=\scriptsize,text=navy]{cm};
\node[font=\tiny,text=muted,below,inner sep=1.5pt] at (1,0) {$1$};
\node[font=\tiny,text=muted,below,inner sep=1.5pt] at (2,0) {$2$};
\node[font=\tiny,text=muted,below,inner sep=1.5pt] at (3,0) {$3$};
\node[font=\tiny,text=muted,below,inner sep=1.5pt] at (4,0) {$4$};
\node[font=\tiny,text=muted,below,inner sep=1.5pt] at (5,0) {$5$};
\node[font=\tiny,text=muted,below,inner sep=1.5pt] at (6,0) {$6$};
\node[font=\tiny,text=muted,left,inner sep=1.5pt] at (0,2) {$2$};
\node[font=\tiny,text=muted,left,inner sep=1.5pt] at (0,4) {$4$};
\node[font=\tiny,text=muted,left,inner sep=1.5pt] at (0,6) {$6$};
\node[font=\tiny,text=muted,left,inner sep=1.5pt] at (0,8) {$8$};
\node[font=\tiny,text=muted,left,inner sep=1.5pt] at (0,10) {$10$};
\node[font=\tiny,text=muted,left,inner sep=1.5pt] at (0,12) {$12$};
\node[font=\tiny,text=muted,left,inner sep=1.5pt] at (0,14) {$14$};
\node[font=\tiny,text=muted,left,inner sep=1.5pt] at (0,16) {$16$};
\node[font=\tiny,text=muted,below left,inner sep=1.5pt] at (0,0) {$0$};
\draw[teal,line width=1.1pt] (0,2)--(1,5)--(2,9)--(3,12)--(4,14)--(5,15)--(6,15);
\fill[navy] (0,2) circle (2.2pt);
\fill[navy] (1,5) circle (2.2pt);
\fill[navy] (2,9) circle (2.2pt);
\fill[navy] (3,12) circle (2.2pt);
\fill[navy] (4,14) circle (2.2pt);
\fill[navy] (5,15) circle (2.2pt);
\fill[navy] (6,15) circle (2.2pt);
\end{tikzpicture}\end{center}\begin{tasks}(1)
  \task ¿Cuánto medía al empezar? ¿Y en la 4.\textsuperscript{a} semana?
  \task ¿En qué semana creció más?
  \task ¿Qué pasa entre la 5.\textsuperscript{a} y la 6.\textsuperscript{a} semana?
\end{tasks}

\ex{2} La gráfica muestra la altura del agua de una bañera. Explica qué pasa en cada tramo.\par\smallskip\begin{center}\begin{tikzpicture}[x=0.33cm,y=0.14cm,baseline=(current bounding box.center),line cap=round,line join=round]
\draw[muted!30,line width=0.3pt] (0,0)--(0,30);
\draw[muted!30,line width=0.3pt] (2,0)--(2,30);
\draw[muted!30,line width=0.3pt] (4,0)--(4,30);
\draw[muted!30,line width=0.3pt] (6,0)--(6,30);
\draw[muted!30,line width=0.3pt] (8,0)--(8,30);
\draw[muted!30,line width=0.3pt] (10,0)--(10,30);
\draw[muted!30,line width=0.3pt] (12,0)--(12,30);
\draw[muted!30,line width=0.3pt] (14,0)--(14,30);
\draw[muted!30,line width=0.3pt] (16,0)--(16,30);
\draw[muted!30,line width=0.3pt] (18,0)--(18,30);
\draw[muted!30,line width=0.3pt] (20,0)--(20,30);
\draw[muted!30,line width=0.3pt] (22,0)--(22,30);
\draw[muted!30,line width=0.3pt] (0,0)--(23,0);
\draw[muted!30,line width=0.3pt] (0,5)--(23,5);
\draw[muted!30,line width=0.3pt] (0,10)--(23,10);
\draw[muted!30,line width=0.3pt] (0,15)--(23,15);
\draw[muted!30,line width=0.3pt] (0,20)--(23,20);
\draw[muted!30,line width=0.3pt] (0,25)--(23,25);
\draw[muted!30,line width=0.3pt] (0,30)--(23,30);
\draw[-latex,muted,line width=0.6pt] (0,0)--(23.6,0) node[right,font=\scriptsize,text=navy]{min};
\draw[-latex,muted,line width=0.6pt] (0,0)--(0,31.5) node[above,font=\scriptsize,text=navy]{cm};
\node[font=\tiny,text=muted,below,inner sep=1.5pt] at (2,0) {$2$};
\node[font=\tiny,text=muted,below,inner sep=1.5pt] at (4,0) {$4$};
\node[font=\tiny,text=muted,below,inner sep=1.5pt] at (6,0) {$6$};
\node[font=\tiny,text=muted,below,inner sep=1.5pt] at (8,0) {$8$};
\node[font=\tiny,text=muted,below,inner sep=1.5pt] at (10,0) {$10$};
\node[font=\tiny,text=muted,below,inner sep=1.5pt] at (12,0) {$12$};
\node[font=\tiny,text=muted,below,inner sep=1.5pt] at (14,0) {$14$};
\node[font=\tiny,text=muted,below,inner sep=1.5pt] at (16,0) {$16$};
\node[font=\tiny,text=muted,below,inner sep=1.5pt] at (18,0) {$18$};
\node[font=\tiny,text=muted,below,inner sep=1.5pt] at (20,0) {$20$};
\node[font=\tiny,text=muted,below,inner sep=1.5pt] at (22,0) {$22$};
\node[font=\tiny,text=muted,left,inner sep=1.5pt] at (0,5) {$5$};
\node[font=\tiny,text=muted,left,inner sep=1.5pt] at (0,10) {$10$};
\node[font=\tiny,text=muted,left,inner sep=1.5pt] at (0,15) {$15$};
\node[font=\tiny,text=muted,left,inner sep=1.5pt] at (0,20) {$20$};
\node[font=\tiny,text=muted,left,inner sep=1.5pt] at (0,25) {$25$};
\node[font=\tiny,text=muted,below left,inner sep=1.5pt] at (0,0) {$0$};
\draw[accent,line width=1.1pt] (0,0)--(5,20)--(6,26)--(16,26)--(17,20)--(22,0);
\end{tikzpicture}\end{center}

\ex{2} ¿Qué gráfica (tiempo, distancia a casa) corresponde a cada historia? Hay una que sobra.\par\smallskip\begin{center}\begin{tabular}{c}\begin{tikzpicture}[x=0.3cm,y=0.22cm,baseline=(current bounding box.center),line cap=round,line join=round]
\draw[muted!30,line width=0.3pt] (0,0)--(0,9);
\draw[muted!30,line width=0.3pt] (2,0)--(2,9);
\draw[muted!30,line width=0.3pt] (4,0)--(4,9);
\draw[muted!30,line width=0.3pt] (6,0)--(6,9);
\draw[muted!30,line width=0.3pt] (0,0)--(6.5,0);
\draw[muted!30,line width=0.3pt] (0,2)--(6.5,2);
\draw[muted!30,line width=0.3pt] (0,4)--(6.5,4);
\draw[muted!30,line width=0.3pt] (0,6)--(6.5,6);
\draw[muted!30,line width=0.3pt] (0,8)--(6.5,8);
\draw[-latex,muted,line width=0.6pt] (0,0)--(7.1,0) node[right,font=\scriptsize,text=navy]{$t$};
\draw[-latex,muted,line width=0.6pt] (0,0)--(0,9.6) node[above,font=\scriptsize,text=navy]{$d$};
\node[font=\tiny,text=muted,below,inner sep=1.5pt] at (2,0) {$2$};
\node[font=\tiny,text=muted,below,inner sep=1.5pt] at (4,0) {$4$};
\node[font=\tiny,text=muted,below,inner sep=1.5pt] at (6,0) {$6$};
\node[font=\tiny,text=muted,left,inner sep=1.5pt] at (0,2) {$2$};
\node[font=\tiny,text=muted,left,inner sep=1.5pt] at (0,4) {$4$};
\node[font=\tiny,text=muted,left,inner sep=1.5pt] at (0,6) {$6$};
\node[font=\tiny,text=muted,left,inner sep=1.5pt] at (0,8) {$8$};
\node[font=\tiny,text=muted,below left,inner sep=1.5pt] at (0,0) {$0$};
\draw[accent,line width=1.1pt] (0,0)--(2,4)--(4,4)--(6,0);
\end{tikzpicture}\\ \textbf{A}\end{tabular}\hspace{0.3cm}\begin{tabular}{c}\begin{tikzpicture}[x=0.3cm,y=0.22cm,baseline=(current bounding box.center),line cap=round,line join=round]
\draw[muted!30,line width=0.3pt] (0,0)--(0,9);
\draw[muted!30,line width=0.3pt] (2,0)--(2,9);
\draw[muted!30,line width=0.3pt] (4,0)--(4,9);
\draw[muted!30,line width=0.3pt] (6,0)--(6,9);
\draw[muted!30,line width=0.3pt] (0,0)--(6.5,0);
\draw[muted!30,line width=0.3pt] (0,2)--(6.5,2);
\draw[muted!30,line width=0.3pt] (0,4)--(6.5,4);
\draw[muted!30,line width=0.3pt] (0,6)--(6.5,6);
\draw[muted!30,line width=0.3pt] (0,8)--(6.5,8);
\draw[-latex,muted,line width=0.6pt] (0,0)--(7.1,0) node[right,font=\scriptsize,text=navy]{$t$};
\draw[-latex,muted,line width=0.6pt] (0,0)--(0,9.6) node[above,font=\scriptsize,text=navy]{$d$};
\node[font=\tiny,text=muted,below,inner sep=1.5pt] at (2,0) {$2$};
\node[font=\tiny,text=muted,below,inner sep=1.5pt] at (4,0) {$4$};
\node[font=\tiny,text=muted,below,inner sep=1.5pt] at (6,0) {$6$};
\node[font=\tiny,text=muted,left,inner sep=1.5pt] at (0,2) {$2$};
\node[font=\tiny,text=muted,left,inner sep=1.5pt] at (0,4) {$4$};
\node[font=\tiny,text=muted,left,inner sep=1.5pt] at (0,6) {$6$};
\node[font=\tiny,text=muted,left,inner sep=1.5pt] at (0,8) {$8$};
\node[font=\tiny,text=muted,below left,inner sep=1.5pt] at (0,0) {$0$};
\draw[accent,line width=1.1pt] (0,0)--(6,8);
\end{tikzpicture}\\ \textbf{B}\end{tabular}\hspace{0.3cm}\begin{tabular}{c}\begin{tikzpicture}[x=0.3cm,y=0.22cm,baseline=(current bounding box.center),line cap=round,line join=round]
\draw[muted!30,line width=0.3pt] (0,0)--(0,9);
\draw[muted!30,line width=0.3pt] (2,0)--(2,9);
\draw[muted!30,line width=0.3pt] (4,0)--(4,9);
\draw[muted!30,line width=0.3pt] (6,0)--(6,9);
\draw[muted!30,line width=0.3pt] (0,0)--(6.5,0);
\draw[muted!30,line width=0.3pt] (0,2)--(6.5,2);
\draw[muted!30,line width=0.3pt] (0,4)--(6.5,4);
\draw[muted!30,line width=0.3pt] (0,6)--(6.5,6);
\draw[muted!30,line width=0.3pt] (0,8)--(6.5,8);
\draw[-latex,muted,line width=0.6pt] (0,0)--(7.1,0) node[right,font=\scriptsize,text=navy]{$t$};
\draw[-latex,muted,line width=0.6pt] (0,0)--(0,9.6) node[above,font=\scriptsize,text=navy]{$d$};
\node[font=\tiny,text=muted,below,inner sep=1.5pt] at (2,0) {$2$};
\node[font=\tiny,text=muted,below,inner sep=1.5pt] at (4,0) {$4$};
\node[font=\tiny,text=muted,below,inner sep=1.5pt] at (6,0) {$6$};
\node[font=\tiny,text=muted,left,inner sep=1.5pt] at (0,2) {$2$};
\node[font=\tiny,text=muted,left,inner sep=1.5pt] at (0,4) {$4$};
\node[font=\tiny,text=muted,left,inner sep=1.5pt] at (0,6) {$6$};
\node[font=\tiny,text=muted,left,inner sep=1.5pt] at (0,8) {$8$};
\node[font=\tiny,text=muted,below left,inner sep=1.5pt] at (0,0) {$0$};
\draw[accent,line width=1.1pt] (0,0)--(2,4)--(4,4)--(6,8);
\end{tikzpicture}\\ \textbf{C}\end{tabular}\hspace{0.3cm}\begin{tabular}{c}\begin{tikzpicture}[x=0.3cm,y=0.22cm,baseline=(current bounding box.center),line cap=round,line join=round]
\draw[muted!30,line width=0.3pt] (0,0)--(0,9);
\draw[muted!30,line width=0.3pt] (2,0)--(2,9);
\draw[muted!30,line width=0.3pt] (4,0)--(4,9);
\draw[muted!30,line width=0.3pt] (6,0)--(6,9);
\draw[muted!30,line width=0.3pt] (0,0)--(6.5,0);
\draw[muted!30,line width=0.3pt] (0,2)--(6.5,2);
\draw[muted!30,line width=0.3pt] (0,4)--(6.5,4);
\draw[muted!30,line width=0.3pt] (0,6)--(6.5,6);
\draw[muted!30,line width=0.3pt] (0,8)--(6.5,8);
\draw[-latex,muted,line width=0.6pt] (0,0)--(7.1,0) node[right,font=\scriptsize,text=navy]{$t$};
\draw[-latex,muted,line width=0.6pt] (0,0)--(0,9.6) node[above,font=\scriptsize,text=navy]{$d$};
\node[font=\tiny,text=muted,below,inner sep=1.5pt] at (2,0) {$2$};
\node[font=\tiny,text=muted,below,inner sep=1.5pt] at (4,0) {$4$};
\node[font=\tiny,text=muted,below,inner sep=1.5pt] at (6,0) {$6$};
\node[font=\tiny,text=muted,left,inner sep=1.5pt] at (0,2) {$2$};
\node[font=\tiny,text=muted,left,inner sep=1.5pt] at (0,4) {$4$};
\node[font=\tiny,text=muted,left,inner sep=1.5pt] at (0,6) {$6$};
\node[font=\tiny,text=muted,left,inner sep=1.5pt] at (0,8) {$8$};
\node[font=\tiny,text=muted,below left,inner sep=1.5pt] at (0,0) {$0$};
\draw[accent,line width=1.1pt] (0,8)--(3,2)--(6,2);
\end{tikzpicture}\\ \textbf{D}\end{tabular}\end{center}\begin{tasks}(1)
  \task Sara va a casa de una amiga sin parar.
  \task Óscar sale de casa, se para a comprar el pan y continúa hasta la escuela.
  \task Lucía va al quiosco, se queda un rato y vuelve a casa.
\end{tasks}

\ex{2} La gráfica muestra el viaje de un tren.\par\smallskip\begin{center}\begin{tikzpicture}[x=0.12cm,y=0.05cm,baseline=(current bounding box.center),line cap=round,line join=round]
\draw[muted!30,line width=0.3pt] (0,0)--(0,100);
\draw[muted!30,line width=0.3pt] (5,0)--(5,100);
\draw[muted!30,line width=0.3pt] (10,0)--(10,100);
\draw[muted!30,line width=0.3pt] (15,0)--(15,100);
\draw[muted!30,line width=0.3pt] (20,0)--(20,100);
\draw[muted!30,line width=0.3pt] (25,0)--(25,100);
\draw[muted!30,line width=0.3pt] (30,0)--(30,100);
\draw[muted!30,line width=0.3pt] (35,0)--(35,100);
\draw[muted!30,line width=0.3pt] (40,0)--(40,100);
\draw[muted!30,line width=0.3pt] (45,0)--(45,100);
\draw[muted!30,line width=0.3pt] (50,0)--(50,100);
\draw[muted!30,line width=0.3pt] (55,0)--(55,100);
\draw[muted!30,line width=0.3pt] (60,0)--(60,100);
\draw[muted!30,line width=0.3pt] (65,0)--(65,100);
\draw[muted!30,line width=0.3pt] (70,0)--(70,100);
\draw[muted!30,line width=0.3pt] (75,0)--(75,100);
\draw[muted!30,line width=0.3pt] (0,0)--(75,0);
\draw[muted!30,line width=0.3pt] (0,10)--(75,10);
\draw[muted!30,line width=0.3pt] (0,20)--(75,20);
\draw[muted!30,line width=0.3pt] (0,30)--(75,30);
\draw[muted!30,line width=0.3pt] (0,40)--(75,40);
\draw[muted!30,line width=0.3pt] (0,50)--(75,50);
\draw[muted!30,line width=0.3pt] (0,60)--(75,60);
\draw[muted!30,line width=0.3pt] (0,70)--(75,70);
\draw[muted!30,line width=0.3pt] (0,80)--(75,80);
\draw[muted!30,line width=0.3pt] (0,90)--(75,90);
\draw[muted!30,line width=0.3pt] (0,100)--(75,100);
\draw[-latex,muted,line width=0.6pt] (0,0)--(76.5,0) node[right,font=\scriptsize,text=navy]{min};
\draw[-latex,muted,line width=0.6pt] (0,0)--(0,103) node[above,font=\scriptsize,text=navy]{km};
\node[font=\tiny,text=muted,below,inner sep=1.5pt] at (10,0) {$10$};
\node[font=\tiny,text=muted,below,inner sep=1.5pt] at (20,0) {$20$};
\node[font=\tiny,text=muted,below,inner sep=1.5pt] at (30,0) {$30$};
\node[font=\tiny,text=muted,below,inner sep=1.5pt] at (40,0) {$40$};
\node[font=\tiny,text=muted,below,inner sep=1.5pt] at (50,0) {$50$};
\node[font=\tiny,text=muted,below,inner sep=1.5pt] at (60,0) {$60$};
\node[font=\tiny,text=muted,below,inner sep=1.5pt] at (70,0) {$70$};
\node[font=\tiny,text=muted,left,inner sep=1.5pt] at (0,20) {$20$};
\node[font=\tiny,text=muted,left,inner sep=1.5pt] at (0,40) {$40$};
\node[font=\tiny,text=muted,left,inner sep=1.5pt] at (0,60) {$60$};
\node[font=\tiny,text=muted,left,inner sep=1.5pt] at (0,80) {$80$};
\node[font=\tiny,text=muted,below left,inner sep=1.5pt] at (0,0) {$0$};
\draw[accent,line width=1.1pt] (0,0)--(20,30)--(25,30)--(45,60)--(50,60)--(70,90);
\end{tikzpicture}\end{center}\begin{tasks}(1)
  \task ¿Cuántas paradas hace y cuánto duran?
  \task ¿Cuántos km hace entre parada y parada? ¿Cuánto tarda?
  \task ¿Qué distancia total recorre y cuánto tiempo tarda?
\end{tasks}

\ex{3} Ana y Pablo hacen la misma ruta en bici. Observa la gráfica y responde.\par\smallskip\begin{center}\begin{tikzpicture}[x=0.9cm,y=0.04cm,baseline=(current bounding box.center),line cap=round,line join=round]
\draw[muted!30,line width=0.3pt] (0,0)--(0,105);
\draw[muted!30,line width=0.3pt] (1,0)--(1,105);
\draw[muted!30,line width=0.3pt] (2,0)--(2,105);
\draw[muted!30,line width=0.3pt] (3,0)--(3,105);
\draw[muted!30,line width=0.3pt] (4,0)--(4,105);
\draw[muted!30,line width=0.3pt] (5,0)--(5,105);
\draw[muted!30,line width=0.3pt] (0,0)--(5.5,0);
\draw[muted!30,line width=0.3pt] (0,10)--(5.5,10);
\draw[muted!30,line width=0.3pt] (0,20)--(5.5,20);
\draw[muted!30,line width=0.3pt] (0,30)--(5.5,30);
\draw[muted!30,line width=0.3pt] (0,40)--(5.5,40);
\draw[muted!30,line width=0.3pt] (0,50)--(5.5,50);
\draw[muted!30,line width=0.3pt] (0,60)--(5.5,60);
\draw[muted!30,line width=0.3pt] (0,70)--(5.5,70);
\draw[muted!30,line width=0.3pt] (0,80)--(5.5,80);
\draw[muted!30,line width=0.3pt] (0,90)--(5.5,90);
\draw[muted!30,line width=0.3pt] (0,100)--(5.5,100);
\draw[-latex,muted,line width=0.6pt] (0,0)--(5.8,0) node[right,font=\scriptsize,text=navy]{h};
\draw[-latex,muted,line width=0.6pt] (0,0)--(0,108) node[above,font=\scriptsize,text=navy]{km};
\node[font=\tiny,text=muted,below,inner sep=1.5pt] at (1,0) {$1$};
\node[font=\tiny,text=muted,below,inner sep=1.5pt] at (2,0) {$2$};
\node[font=\tiny,text=muted,below,inner sep=1.5pt] at (3,0) {$3$};
\node[font=\tiny,text=muted,below,inner sep=1.5pt] at (4,0) {$4$};
\node[font=\tiny,text=muted,below,inner sep=1.5pt] at (5,0) {$5$};
\node[font=\tiny,text=muted,left,inner sep=1.5pt] at (0,20) {$20$};
\node[font=\tiny,text=muted,left,inner sep=1.5pt] at (0,40) {$40$};
\node[font=\tiny,text=muted,left,inner sep=1.5pt] at (0,60) {$60$};
\node[font=\tiny,text=muted,left,inner sep=1.5pt] at (0,80) {$80$};
\node[font=\tiny,text=muted,left,inner sep=1.5pt] at (0,100) {$100$};
\node[font=\tiny,text=muted,below left,inner sep=1.5pt] at (0,0) {$0$};
\draw[accent,line width=1.1pt] (0,0)--(5,75);
\draw[teal,line width=1.1pt] (1,0)--(5,100);
\fill[navy] (2.5,37.5) circle (2.2pt);
\node[font=\scriptsize,right,text=accent] at (5,75) {Ana};
\node[font=\scriptsize,left,text=teal] at (4.6,100) {Pablo};
\end{tikzpicture}\end{center}\begin{tasks}(1)
  \task ¿Quién sale primero? ¿Cuánto rato antes?
  \task ¿Cuántos km hace cada uno en una hora? ¿Quién va más deprisa?
  \task ¿Dónde y cuándo alcanza Pablo a Ana?
\end{tasks}

\ex{3} Marcos sale de casa, camina 1\,km en 20 minutos hasta la biblioteca, se queda allí media hora y vuelve a casa en 10 minutos. Dibuja la gráfica (tiempo en minutos, distancia a casa en km). ¿En qué tramo va más deprisa?


% =====================================================
\seccio{5}{Función de proporcionalidad directa: $y=m\cdot x$}

\ex{1} ¿Cuáles de estas funciones son de proporcionalidad directa? Escribe su pendiente $m$.\begin{tasks}(6)
  \task $y=4x$
  \task $y=x+4$
  \task $y=-\tfrac{x}{2}$
  \task $y=4$
  \task $y=0{,}6x$
  \task $y=-3x$
\end{tasks}\recorda{Una función de proporcionalidad directa es de la forma $y=m\cdot x$. Su gráfica es una recta que pasa por el origen. La pendiente $m$ es lo que aumenta $y$ cuando $x$ aumenta 1.}

\ex{1} Completa la tabla de $y=3x$ y dibuja su gráfica.\par\smallskip\begin{center}$\begin{array}{|c|C{0.75cm}C{0.75cm}C{0.75cm}C{0.75cm}C{0.75cm}|}\hline x&-2&-1&0&1&2\\\hline y&\phantom{00}&\phantom{00}&\phantom{00}&\phantom{00}&\phantom{00}\\\hline\end{array}$\end{center}

\ex{1} Dibuja en los mismos ejes $y=-2x$ e $y=\dfrac{x}{2}$. ¿Cuál es creciente? ¿Cuál está más inclinada?

\ex{2} Una recta pasa por el origen y por el punto $(3,\ -6)$. Escribe su fórmula y comprueba si los puntos $(-4,\ 8)$ y $(5,\ -12)$ son de la recta.

\ex{2} Comprueba que esta tabla es de proporcionalidad directa, escribe su fórmula y dibuja la gráfica. ¿Cuánto vale $y$ si $x=10$?\par\smallskip\begin{center}$\begin{array}{|c|cccc|}\hline x&2&4&6&8\\\hline y&3&6&9&12\\\hline\end{array}$\end{center}

\ex{2} Ordena estas rectas de más inclinada a menos y di cuáles son decrecientes: \ $a$: $y=5x$, \ $b$: $y=-\frac12x$, \ $c$: $y=\frac13x$, \ $d$: $y=-2x$.

\ex{3} Un día, 1\,€ vale 1,10 dólares. Escribe la fórmula que pasa de euros ($x$) a dólares ($y$) y dibuja su gráfica. ¿Cuántos dólares son 20\,€? ¿Y 35\,€? ¿Cuántos euros son 55\,\$?

\ex{3} Un coche va a 90\,km/h siempre a la misma velocidad. Escribe la fórmula de la distancia $d$ que recorre en $t$ horas, dibuja su gráfica y calcula qué distancia hace en 3 horas y cuánto tarda en hacer 225\,km.


% =====================================================
\seccio{6}{Función lineal: $y=m\cdot x+n$}

\ex{1} Escribe la pendiente $m$ y la ordenada en el origen $n$ de cada función:\begin{tasks}(6)
  \task $y=3x-4$
  \task $y=-x+5$
  \task $y=\tfrac{x}{2}+1$
  \task $y=-2x$
  \task $y=7$
  \task $y=6-4x$
\end{tasks}\recorda{En $y=mx+n$, $m$ es la \textbf{pendiente} (lo que cambia $y$ cuando $x$ aumenta 1) y $n$ es la \textbf{ordenada en el origen}: la recta corta el eje $Y$ en $(0,\ n)$.}

\ex{1} Completa la tabla de $y=2x-1$ y dibuja su gráfica. ¿Dónde corta el eje $Y$?\par\smallskip\begin{center}$\begin{array}{|c|C{0.75cm}C{0.75cm}C{0.75cm}C{0.75cm}C{0.75cm}|}\hline x&-1&0&1&2&3\\\hline y&\phantom{00}&\phantom{00}&\phantom{00}&\phantom{00}&\phantom{00}\\\hline\end{array}$\end{center}

\ex{1} ¿Cuáles de estos puntos son de la recta $y=-2x+3$? \ $(2,\ 1)$, \ $(-1,\ 7)$, \ $(3,\ -2)$, \ $(0,\ 3)$.

\ex{2} Dibuja en los mismos ejes $y=x+2$ e $y=-x+4$. ¿En qué punto se cortan? Compruébalo con las fórmulas.

\ex{2} Escribe la fórmula de cada una de las rectas del dibujo.\par\smallskip\begin{center}\begin{tikzpicture}[x=0.4cm,y=0.4cm,baseline=(current bounding box.center),line cap=round,line join=round]
\draw[muted!30,line width=0.3pt] (-5,-5)--(-5,5);
\draw[muted!30,line width=0.3pt] (-4,-5)--(-4,5);
\draw[muted!30,line width=0.3pt] (-3,-5)--(-3,5);
\draw[muted!30,line width=0.3pt] (-2,-5)--(-2,5);
\draw[muted!30,line width=0.3pt] (-1,-5)--(-1,5);
\draw[muted!30,line width=0.3pt] (0,-5)--(0,5);
\draw[muted!30,line width=0.3pt] (1,-5)--(1,5);
\draw[muted!30,line width=0.3pt] (2,-5)--(2,5);
\draw[muted!30,line width=0.3pt] (3,-5)--(3,5);
\draw[muted!30,line width=0.3pt] (4,-5)--(4,5);
\draw[muted!30,line width=0.3pt] (5,-5)--(5,5);
\draw[muted!30,line width=0.3pt] (-5,-5)--(5,-5);
\draw[muted!30,line width=0.3pt] (-5,-4)--(5,-4);
\draw[muted!30,line width=0.3pt] (-5,-3)--(5,-3);
\draw[muted!30,line width=0.3pt] (-5,-2)--(5,-2);
\draw[muted!30,line width=0.3pt] (-5,-1)--(5,-1);
\draw[muted!30,line width=0.3pt] (-5,0)--(5,0);
\draw[muted!30,line width=0.3pt] (-5,1)--(5,1);
\draw[muted!30,line width=0.3pt] (-5,2)--(5,2);
\draw[muted!30,line width=0.3pt] (-5,3)--(5,3);
\draw[muted!30,line width=0.3pt] (-5,4)--(5,4);
\draw[muted!30,line width=0.3pt] (-5,5)--(5,5);
\draw[-latex,muted,line width=0.6pt] (-5,0)--(5.3,0) node[right,font=\scriptsize,text=navy]{x};
\draw[-latex,muted,line width=0.6pt] (0,-5)--(0,5.3) node[above,font=\scriptsize,text=navy]{y};
\node[font=\tiny,text=muted,below,inner sep=1.5pt] at (-4,0) {$-4$};
\node[font=\tiny,text=muted,below,inner sep=1.5pt] at (-3,0) {$-3$};
\node[font=\tiny,text=muted,below,inner sep=1.5pt] at (-2,0) {$-2$};
\node[font=\tiny,text=muted,below,inner sep=1.5pt] at (-1,0) {$-1$};
\node[font=\tiny,text=muted,below,inner sep=1.5pt] at (1,0) {$1$};
\node[font=\tiny,text=muted,below,inner sep=1.5pt] at (2,0) {$2$};
\node[font=\tiny,text=muted,below,inner sep=1.5pt] at (3,0) {$3$};
\node[font=\tiny,text=muted,below,inner sep=1.5pt] at (4,0) {$4$};
\node[font=\tiny,text=muted,left,inner sep=1.5pt] at (0,-4) {$-4$};
\node[font=\tiny,text=muted,left,inner sep=1.5pt] at (0,-3) {$-3$};
\node[font=\tiny,text=muted,left,inner sep=1.5pt] at (0,-2) {$-2$};
\node[font=\tiny,text=muted,left,inner sep=1.5pt] at (0,-1) {$-1$};
\node[font=\tiny,text=muted,left,inner sep=1.5pt] at (0,1) {$1$};
\node[font=\tiny,text=muted,left,inner sep=1.5pt] at (0,2) {$2$};
\node[font=\tiny,text=muted,left,inner sep=1.5pt] at (0,3) {$3$};
\node[font=\tiny,text=muted,left,inner sep=1.5pt] at (0,4) {$4$};
\node[font=\tiny,text=muted,below left,inner sep=1.5pt] at (0,0) {$0$};
\draw[accent,line width=1.1pt] (-2.000,-5.000)--(3.000,5.000);
\draw[teal,line width=1.1pt] (-2.000,5.000)--(5.000,-2.000);
\draw[navy,line width=1.1pt] (-5.000,-1.500)--(5.000,3.500);
\draw[violet,line width=1.1pt] (-5.000,-3.000)--(5.000,-3.000);
\node[font=\scriptsize,right,text=accent] at (2.8,4.6) {$a$};
\node[font=\scriptsize,left,text=teal] at (-1.8,4.6) {$b$};
\node[font=\scriptsize,above,text=navy] at (4.4,3.2) {$c$};
\node[font=\scriptsize,above,text=violet] at (-4.5,-3) {$d$};
\end{tikzpicture}\end{center}

\ex{2} Observa estas funciones: \ $y=3x+1$, \ $y=-x$, \ $y=3x-4$, \ $y=5$, \ $y=-x+2$, \ $y=2$. ¿Cuáles son paralelas? ¿Cuáles pasan por el origen? ¿Cuáles son constantes?

\ex{2} Halla la fórmula de la función lineal de esta tabla:\par\smallskip\begin{center}$\begin{array}{|c|cccc|}\hline x&0&1&2&3\\\hline y&4&1&-2&-5\\\hline\end{array}$\end{center}

\ex{3} Halla la fórmula de la recta que pasa por:\begin{tasks}(2)
  \task $(0,\ 3)$ y $(2,\ 7)$
  \task $(1,\ -1)$ y $(3,\ 3)$
\end{tasks}

\ex{3} Responde:\begin{tasks}(1)
  \task ¿Qué valor tiene que tener $k$ para que el punto $(2,\ k)$ sea de la recta $y=-3x+5$?
  \task ¿Qué valor tiene que tener $a$ para que el punto $(a,\ 11)$ sea de la recta $y=2x+5$?
\end{tasks}


% =====================================================
\seccio{7}{Problemas con funciones lineales}

\ex{1} Un taxi cobra 3,50\,€ de bajada de bandera y 1,50\,€ por km. Escribe la fórmula del precio, dibuja su gráfica y calcula cuánto cuesta un viaje de 8\,km. ¿Cuántos km hemos hecho si pagamos 26\,€?\recorda{En muchos problemas hay una \textbf{parte fija} $n$ y una cantidad $m$ que se suma (o se resta) por cada unidad: $y=mx+n$.}

\ex{1} Una vela mide 24\,cm y se consume 3\,cm cada hora. Escribe la fórmula de la altura $y$ al cabo de $x$ horas y dibuja su gráfica. ¿Cuánto medirá al cabo de 5\,h? ¿Cuándo se acabará?

\ex{1} Una tarifa de teléfono cuesta 6\,€ al mes más 0,08\,€ por minuto. ¿Cuánto pagamos si hablamos 50 minutos? ¿Cuántos minutos hemos hablado si pagamos 9\,€?

\ex{2} Un depósito tiene 100\,L de agua y un grifo echa 25\,L cada minuto. Escribe la fórmula. ¿Cuántos litros habrá al cabo de 6 minutos? Si caben 600\,L, ¿cuándo se llenará?

\ex{2} Un fontanero cobra según esta tabla. Halla la fórmula y calcula cuánto cobrará por 6 horas. ¿Cuántas horas ha trabajado si cobra 300\,€?\par\smallskip\begin{center}$\begin{array}{|c|ccc|}\hline \text{Horas}&1&2&3\\\hline \text{Precio (€)}&55&90&125\\\hline\end{array}$\end{center}

\ex{2} Para pasar de grados Celsius ($x$) a grados Fahrenheit ($y$) se usa la fórmula $y=1{,}8x+32$. ¿A cuántos $^\circ$F equivalen 0\,$^\circ$C, 25\,$^\circ$C y 100\,$^\circ$C? ¿A cuántos $^\circ$C equivalen 50\,$^\circ$F?

\ex{2} Un aparcamiento cobra 3\,€ por entrar y 2\,€ por cada hora. Escribe la fórmula y haz una tabla para 1, 2, 3, 4 y 5 horas. ¿Cuántas horas hemos dejado allí el coche si pagamos 19\,€?

\ex{3} Para alquilar una furgoneta, la empresa A cobra 30\,€ más 0,25\,€ por km y la empresa B, 50\,€ más 0,15\,€ por km. ¿Cuál conviene si hacemos 150\,km? ¿Para cuántos km cuestan lo mismo? Haz la gráfica de las dos.

\ex{3} El gimnasio A cobra 45\,€ de matrícula y 20\,€ al mes; el gimnasio B no tiene matrícula y cobra 35\,€ al mes. Haz la gráfica de las dos tarifas. ¿Cuándo cuestan lo mismo? ¿Cuál conviene si quieres ir 6 meses?

\vfill
\begin{center}\color{muted}\small\faCheckCircle[regular]\ \ ¡Ánimo! Si te cuesta algún ejercicio, repasa los apuntes interactivos de la unidad.\end{center}

\end{document}
